\documentclass[12pt]{article}
\usepackage[utf8]{inputenc}
\usepackage[T1]{fontenc}
\usepackage{hyperref}
\title{Paic\^i Language Resources for K-12 Curriculum Development}
\author{Compiled by Claude Sonnet 5.5 (claude-sonnet-5-5) for Andr\'as Kornai}
\date{}
\begin{document}
\maketitle
\section{Introduction}
Paic\^i is a Kanak language of New Caledonia with contrastive tone. Searches found a classic dictionary and some academic work, but no school materials.
\section{Primary Educational Resources}
No primers or textbooks were found. A 2023 language snapshot mentions that a pedagogical grammar is being prepared by one author; no title or URL was seen.
\section{Linguistic and Reference Materials}
\begin{itemize}
\item Rivierre, J.-C., \emph{Dictionnaire pai\^ci--fran\c{c}ais} (1983, SELAF, Paris), mentioned in the search summary; no direct URL for the book was seen.
\item Language snapshot on Pai\^c\^i (Language Documentation \& Description journal): \url{https://lddjournal.org/article/id/1017}
\item Pai\^c\^i language (Wikipedia): \url{https://en.wikipedia.org/wiki/Paic%C3%AE_language}
\item Pai\^c\^i in the Austronesian typological database: \url{https://lpan.eva.mpg.de/austronesian/language.php?id=1479}
\end{itemize}
\section{Religious and Community Resources}
None found.
\section{Digital and Audio Resources}
A project of the Acad\'emie des Langues Kanak with Princeton University is preparing Kanak dictionaries, including Pai\^c\^i, for public access: \url{https://outremers360.com/bassin-pacifique-appli/nouvelle-caledonie-des-dictionnaires-en-langues-kanak-bientot-accessibles-au-grand-public-grace-a-un-partenariat-avec-luniversite-de-princeton}
\section{Curriculum Development Context}
Little pedagogical material was found; Rivierre's dictionary and the forthcoming digital dictionaries are the main foundations.
\end{document}